\ifx\coursemacrosfileloaded\undefined\gdef\coursemacrosfileloaded{1}\else\endinput\fi
\usepackage{amsthm,amsfonts,amsmath,amssymb}

\usepackage{xcolor}

\usepackage{hyperref}
\hypersetup{
    colorlinks,
    linkcolor={red!50!black},
    citecolor={blue!50!black},
    urlcolor={blue!80!black}
}

\usepackage[most]{tcolorbox}
\usepackage{listings}
\lstdefinestyle{pseudocode}{
  basicstyle=\ttfamily,
  columns=fullflexible,
  mathescape=true,
  frame=single,
  keepspaces=true,
  numbers=none
}


\usepackage{float}

\usepackage{braket}
\usepackage{graphicx,epsfig}
\usepackage{algorithm}
\usepackage{algpseudocode}

\usepackage{tikz}
\usepgflibrary{arrows}
\usetikzlibrary{fit}
\usetikzlibrary{decorations.pathreplacing}
\usetikzlibrary{decorations.markings}
%% PATTERN
\usetikzlibrary{patterns}


%\DeclareMathOperator{\i}{i}
%\DeclareMathOperator{\id}{\boldsymbol{1}}




%\newcommand{\Tr}{\text{Tr}}
%\renewcommand\bra[1]{{\langle{#1}|}}
%\renewcommand\ket[1]{{|{#1}\rangle}}
%\newcommand*{\I}{\imath}%
%\usepackage{graphicx,epsfig}
%\graphicspath{c:/dropbox/teaching/courseNotes/figures/}
%\lstset{language=Python}
%\newtheorem{theorem}{Theorem}

\newcommand{\myop}[1]{\operatorname{#1}}
\newcommand{\inorm}[2]{\left\lVert{#1}\right\rVert_{#2}}
\newcommand{\infnorm}[1]{\inorm{#1}{}}
\newcommand\myOmega{\amalg}
\newcommand{\glen}[1]{\myop{len}_{#1}}
\newcommand{\Rp}{\R_{+}}
